\documentclass[10pt,a4paper]{article}
\usepackage{amsmath,amssymb,amsthm,graphicx}
% Card-local minimal preamble (journal unit ahl-197-on-the-existence-of-logarithmic-and-orbifold-jet-2024, P2-ahl16 precedent).
%% Verbatim from the pinned source: the author theorem environments and macro block,
%% so every macro and environment of the counted spans is available.
% Omitted and why: the documentclass line (the journal class sigma is not on the host;
%% the assembler supplies the class) and page-geometry/font blocks (they fight the
%% forced xelatex engine).
\usepackage{amsmath,amssymb,amsthm,mathtools}

%% \usepackage{physics} is deliberately NOT loaded (L60449 incident 2026-10-15):
%% the span's physics-trigger macros \abs and \Tr are both defined by the author
%% block itself (\DeclarePairedDelimiter{\abs} at pinned source L24, \Tr below),
%% and physics's built-ins collide with them (\Tr, \rank) while its starred
%% delimiter machinery (\opbraces) chokes on the span's multi-line \abs* argument
%% that the author's mathtools version accepts.  The author's own paired
%% delimiter definitions are restored verbatim below (keep-pattern regex gap,
%% fixed 2026-10-15: DeclarePairedDelimiter lines are now kept).
%\DeclareMathOperator{\Im}{Im}
\DeclareMathOperator{\lcm}{lcm}
\DeclareMathOperator{\Arg}{Arg}
\DeclareMathOperator{\Hom}{Hom}
\DeclareMathOperator{\SM}{SM}
\DeclareMathOperator{\Herm}{Herm}
\DeclareMathOperator{\Tr}{Tr}
\DeclareMathOperator{\Pic}{Pic}
\DeclareMathOperator{\Proj}{Proj}
\DeclareMathOperator{\Vol}{Vol}
\DeclareMathOperator{\Span}{Span}
\DeclareMathOperator{\rank}{rank}
\DeclareMathOperator{\diff}{d}
\DeclareMathOperator{\dbar}{\bar\partial}
\DeclareMathOperator{\ddbar}{\partial\bar\partial}
%% Author's own paired delimiters, verbatim from pinned source lines 24-25
%% (the keep-pattern regex lacked DeclarePairedDelimiter; fixed 2026-10-15).
\DeclarePairedDelimiter{\abs}{\lvert}{\rvert}
\DeclarePairedDelimiter{\Abs}{\lVert}{\rVert}
\newenvironment{equation'}{
\begin{equation}
\tag{\theequation\ensuremath{'}}}{
\end{equation}
}
\newenvironment{equation_}[1]{
\begin{equation}
\tag{\ensuremath{\theequation_{#1}}}}{
\end{equation}
}
\newenvironment{eqlist}{
\begingroup \leftskip\parindent \rightskip\parindent \noindent
\begin{minipage}{\linewidth-2\parindent}} {
\end{minipage}
\endgroup}


%\let\Im\relax




%\let\Im\relax
\newcommand*\rS{\mathrm{S}}
\newcommand*\Id{\mathrm{Id}}
\newcommand*\Sing{\mathrm{Sing}}
\newcommand*\orb{\mathrm{orb}}
\newcommand*\FS{\mathrm{FS}}
\newcommand*\loc{\mathrm{loc}}
\newcommand{\rma}{\mathrm{a}}
\newcommand{\rmb}{\mathrm{b}}
\newcommand*\cI{\mathcal{I}}
\newcommand*\cJ{\mathcal{J}}
\newcommand*\cO{\mathcal{O}}
\newcommand*\bC{\mathbb{C}}
\newcommand*\bD{\mathbb{D}}
\newcommand*\bE{\mathbb{E}}
\newcommand*\bG{\mathbb{G}}
\newcommand*\bN{\mathbb{N}}
\newcommand*\bP{\mathbb{P}}
\newcommand*\bQ{\mathbb{Q}}
\newcommand*\bR{\mathbb{R}}
\newcommand*\bS{\mathbb{S}}
\newcommand*\bZ{\mathbb{Z}}
\newcommand*\bB{\mathbb{B}}
\newcommand*\bOne{{\mathchoice {\rm 1\mskip-4mu l} {\rm 1\mskip-4mu l} {\rm 1\mskip-4.5mu l}{\rm 1\mskip-5mu l}}}
\newcommand\bDel{\mbox{$\Delta$\newline \scalebox{.6}{\raisebox{3pt}{$\mkern-22mu\not{}\mkern12mu$}}}}
%\newcommand\bDel{\ooalign{$\Delta$\\ \scalebox{.6}{\raisebox{3pt}{\,$\not{}$}}}}
\newcommand\bDelta{\mathchoice{\bDel}{\bDel}{\scalebox{.7}{\bDel}}{\scalebox{.5}{\bDel}}}
%\newcommand\bDelta{\mbox{$\Delta\mkern-2.5mu\not{}$}}
\let\leq\leqslant
\let\geq\geqslant
\newcommand*\bydef{\coloneqq}
\newcommand*{\mk}{\mkern -1mu}
\newcommand*{\Mk}{\mkern -2mu}
\newcommand*{\mK}{\mkern 1mu}
\newcommand*{\MK}{\mkern 2mu}
\newcommand*{\romanenumi}{\renewcommand*{\theenumi}{\roman{enumi}}}
\newcommand*{\Romanenumi}{\renewcommand*{\theenumi}{\Roman{enumi}}}
\newcommand*{\alphenumi}{\renewcommand*{\theenumi}{\alph{enumi}}}
\newcommand*{\Alphenumi}{\renewcommand*{\theenumi}{\Alph{enumi}}}
\let\oldtilde\tilde
\renewcommand*{\tilde}[1]{\mathchoice{\widetilde{#1}}{\widetilde{#1}}{\oldtilde{#1}}{\oldtilde{#1}}}
%\let\oldhat\hat
%\renewcommand*{\hat}[1]{\mathchoice{\widehat{#1}}{\widehat{#1}}{\oldhat{#1}}{\oldhat{#1}}}
\let\oldexists\exists
\renewcommand*{\exists}{\mathrel{\oldexists}}
\let\oldforall\forall
\renewcommand*{\forall}{\mathrel{\oldforall}}
\newtheorem{cdrthm}{Theorem}[section]
\newtheorem{theo}[cdrthm]{Theorem}
\newtheorem{prop}[cdrthm]{Proposition}
\newtheorem{lemm}[cdrthm]{Lemma}
\newtheorem{coro}[cdrthm]{Corollary}
\newtheorem{defi}[cdrthm]{Definition}
\newtheorem{rema}[cdrthm]{Remark}
\newtheorem{conj}[cdrthm]{Conjecture}
\numberwithin{equation}{section}

%% compile hygiene: sub-50pt overfull reports are cosmetic;
\hfuzz=50pt
\usepackage{cleveref}
\crefname{originalreference}{original reference}{original references}
\Crefname{originalreference}{Original reference}{Original references}
\makeatletter
\newlabel{theo:orbi_curv_estimate}{{1.10}{source}{}{sourcepointer-1}{}}
\newlabel{theo:orbi_curv_estimate@cref}{{[originalreference][0][]1.10}{[1][0][]source}{}{}{}}
\makeatother
\title{Hermitian metrics on orbifold vector bundles with prescribed curvature dominance}
\author{Campana, Frederic; Darondeau, Lionel; Demailly, Jean-Pierre; Rousseau, Erwan}
\date{doi:10.5802/ahl.197; CC BY 4.0}
\begin{document}
\maketitle
\setcounter{section}{7}
\setcounter{equation}{0}
\setcounter{cdrthm}{0}
\noindent\textit{Editorial note.} All delivered mathematical text is the author's own, copied from the pinned journal source file sources/ahl-197-on-the-existence-of-logarithmic-and-orbifold-jet-2024/pinned-original-source.tex (On the existence of logarithmic and orbifold jet differentials); no statement or proof step was rewritten, reordered or removed.\par\medskip
\noindent\textit{Editorial note.} The counted proof is the author's complete argument: Statement at source lines 1903-1917, proof at lines 1918-2279 (single proof block, gate span-guard 4/4). The counted span carries zero \textbackslash{}cite and no bibliography need. Printed number 7.1 in Section 7 (se:curvature\_tangent, sse:orbi\_metric). Batch-3 rotation band first card (duty split 2026-10-15 19:0x).\par\medskip
\noindent\textit{Editorial note.} The label prop:orbi\_metric names the counted statement, and the proof environment immediately following it in the source is delivered verbatim as its proof.\par\medskip
\noindent\textit{Editorial note.} The counted statement is carried by the author's own prop environment in the article source.\par\medskip
\noindent\textit{Editorial note.} Journal-unit assembly: card-local preamble (author theorem environments and macros verbatim), article class, no bibliography in the counted span (the span carries no \textbackslash{}cite commands); references leaving the span are declared as navigation stubs carrying the article's printed numbers.\par\medskip
\noindent\textit{Editorial scope.} the counted statement and its own proof environment only\par\medskip
\ifdefined\hypertarget\hypertarget{sourcepointer-1}{}\fi
\noindent\textit{Source reference.} Original reference 1.10: sources/ahl-197-on-the-existence-of-logarithmic-and-orbifold-jet-2024/pinned-original-source.tex, source lines 314--314. Navigation pointer invoked at line(s) 2278: the counted statement is phrased in notation from that statement of the same article, which is not retained in this card, so the reference is declared as a pointer carrying the number the article itself prints for it (the section and the place of the statement inside that section, derived by the census v3.1 rule for the shared theorem counter); the cited statement is an allowed input and none of it is reproduced inside the span.\par\medskip
\begin{prop}\label{prop:orbi_metric}
Let $X$ be a projective variety, $A$ an ample line bundle, and $(X,V,D)$ an orbifold directed structure where
$D=\sum_{1\,\leq\,j\,\leq\,N}(1-\tfrac{1}{\rho_j})\Delta_j$ is a normal crossing divisor transverse to $V$ in $X$. Let $d_j$ be the infimum of numbers
$\lambda\in\bR_+$ such that $\lambda A-\Delta_j$ is~nef, and $\gamma_V$ be the infimum of numbers $\gamma\geq 0$ such that $\gamma\,
\Theta_{A,h_A}\otimes\Id_V-\Theta_{V,h_V}\geq_G0$ for suitable smooth hermitian metrics $h_V$ on $V$. Then for every
$\gamma>\gamma_{V,D}\bydef\max(\max_j(d_j/\rho_j),\gamma_V)$, the~orbifold vector bundle $V\langle D\rangle$ possesses a hermitian metric $h_{V\langle D\rangle,\gamma,\varepsilon}$ such that
\begin{enumerate}
\item
$h_{V\langle D\rangle,\gamma,\varepsilon}$ is smooth on $X\setminus\abs{D},$
\item
$h_{V\langle D\rangle,\gamma,\varepsilon}$ has the appropriate orbifold singularities along $D,$
\item we have $\gamma\,\Theta_{A,h_A}\otimes\Id-
\Theta_{V\langle D\rangle,h_{V\langle D\rangle,\gamma,\varepsilon}}\geq_G0$ on $X\setminus\abs{D}$.
\end{enumerate}
\end{prop}

\begin{proof}
Let $h_A$ be a metric on $A$ such that $\Theta_{A,h_A}>0$, written locally as $h_A=e^{-\psi}$, and take
\text{$\gamma>\max(\max_j(d_j/\rho_j),\gamma_V)$}. Consider the tautological sections
$\sigma_j\in H^0(X,\cO_X(\Delta_j))$ defining $\Delta_j=\sigma_j^{-1}(0)$, and let $h_V$, $h_j$ be smooth hermitian metrics on
$V$ and $\cO_X(\Delta_j)$ such that
\begin{align}
\stepcounter{equation}
\tag{$\theequation_0$} &
\gamma\,\Theta_{A,h_A}\otimes\Id_V-\Theta_{V,h_V}>_G0,
\\
\tag{$\theequation_j$}\label{6.2_j}
&
\gamma\,\Theta_{A,h_A}-\tfrac{1}{\rho_j}\Theta_{\cO_X(\Delta_j),h_j}>0,\quad
\forall j=1,\dotsc,N,
\end{align}
as is possible by our choice of the constants $d_j$ and $\gamma$. Finally, denote by $\nabla_j$ the associated Chern connection on
$\cO_X(\Delta_j)$. If we write $h_j=e^{-\varphi_j}$ in some local trivialization, then $\nabla_j\sigma_j=\nabla_j^{1,0}\sigma_j=
\partial\sigma_j-\sigma_j\partial\varphi_j$. Take $\omega_A=\Theta_{A,h_A}$ as the K\"ahler metric on $X$. We have
\[
\imath\ddbar\abs{\sigma_j}_{h_j}^{2/\rho_j}=
\tfrac{1}{\rho_j^2}\,\abs{\sigma_j}_{h_j}^{-2+2/\rho_j}\, i\left\langle\nabla_j\sigma_j,\nabla_j\sigma_j\right\rangle_{h_j} -\tfrac{1}{\rho_j}\abs{\sigma_j}_{h_j}^{2/\rho_j}\,\imath\ddbar\varphi_j,
\]
hence there exists $\delta>0$ small such that the metric
$h_{A,\delta}=h_A\exp(-\delta\sum_j\abs{\sigma_j}_{h_j}^{2/\rho_j})$ of weight $\psi_\delta=\psi+\delta\sum_j\abs{\sigma_j}_{h_j}^{2/\rho_j}$ satisfies
\begin{multline*}
\imath\ddbar\psi_\delta(\xi,\xi) =
\\
\abs{\xi}^2_{\omega_A} +
\delta\imath\ddbar
\sum_j\abs{\sigma_j}_{h_j}^{2/\rho_j}(\xi,\xi)
\ge (1-C\delta)\abs{\xi}^2_{\omega_A} +
\delta\sum_j\tfrac{1}{\rho_j^2}\abs{\sigma_j}_{h_j}^{-2+2/\rho_j}
\left|\nabla_j\sigma_j(\xi)\right|_{h_j}^2.
\end{multline*}
We can consider $\omega_{A,\delta}=\Theta_{A,h_{A,\delta}}=\imath\ddbar\psi_\delta$ as an orbifold K\"ahler metric, that is ``smooth'' from the point of view of the orbifold structure. Let us explain the more precise meaning of this ``orbifold smoothness'' assumption. In fact, there exists a ramified cover $g_Y\colon Y\to X$ such that $g^*\sigma_j=w_j^{m_j}$ for some local coordinate $w_j$ on $Y$, with arbitrary high multiplicity $m_j\in\bN^*$ along $\smash{g_Y^{-1}}(\Delta_j)=\{w_j=0\}$. Then $g_Y^*h_{A,\delta}=g_Y^*h_A\exp(-\delta\sum_j \smash{\abs{w_j}^{2m_j/\rho_j})}$ can be taken in any regularity class $C^p$, $p\in \bN^*$, by taking $m_j\geq p\,\rho_j$. Therefore, by pulling-back our calculations to $Y$, we would actually get forms of high regularity on $Y$. Of course, if we compute an integral over $X$, pulling-back forms to $Y$ multiplies the integral by the degree of $g_Y$, and it suffices to divide by that degree to recover the integral over $X$. For $\delta>0$ sufficiently small, our positivity conditions~\eqref{6.2_j} can be turned into the stronger form
\begin{align}
\stepcounter{equation}
\tag{$\theequation_0$}
\gamma\,\imath\ddbar\psi_\delta(\xi,\xi)\,\abs{u}^2-\widetilde\Theta_{V,h_V}(\xi\otimes u) &\ge c\left(\abs{\xi}_{\omega_A}^2+\sum_j\abs{\sigma_j}_{h_j}^{-2+2/\rho_j}\,\left|\nabla_j\sigma_j(\xi)\right|_{h_j}^2
\right),
\\
\tag{$\theequation_j$}\label{6.3_j}
\gamma\,\imath\ddbar\psi_\delta(\xi,\xi)-\tfrac{1}{\rho_j}\,\imath\ddbar\varphi_j(\xi,\xi) &\ge c\left(\abs{\xi}_{\omega_A}^2+\sum_j\abs{\sigma_j}_{h_j}^{-2+2/\rho_j}\,\left|\nabla_j\sigma_j(\xi)\right|_{h_j}^2
\right),
\end{align}
for some constant $c>0$ and all $\xi\in T_X$, $u\in V$ (observe that the right hand side can in fact be seen as a positive definite hermitian form with respect to the orbifold coordinates, we just exploit the fact that
$A$ remains ample when viewed as a line bundle on the orbifold structure). We are going to estimate the curvature of the orbifold metric
$h_{V\langle D\rangle,\varepsilon}$ on $V\langle D\rangle$ defined by
\begin{equation}\label{eq:orbi_metric}
\Abs{u}_{h_{V\langle D\rangle,\varepsilon}}^2 =
\abs{u}_{h_V}^2+\sum_j\varepsilon_j\,
\abs{\sigma_j}_{h_j}^{-2(1-1/\rho_j)}\,\left|\nabla_j\sigma_j(u)\right|_{h_j}^2,
\quad \varepsilon_j\ll 1.
\end{equation}
Again, this metric can be seen as orbifold smooth (in the sense that the metric $g_Y^*h_{V\langle D\rangle,\varepsilon}$ on $g_Y^*(V\langle D\rangle)$ may be taken of arbitrary high regularity; in case $\rho_j=\infty$, it is actually a smooth metric on the logarithmic bundle). Since
\[
\imath\ddbar\Abs{u}_{h_{V\langle D\rangle},\varepsilon}^2=
\imath\langle\nabla u,\nabla u\rangle_{h_{V\langle D\rangle},\varepsilon}-2\pi\,
\left\langle\Theta_{V\langle D\rangle,h_{V\langle D\rangle,\varepsilon}}(u),u
\right\rangle_{h_{V\langle D\rangle,\varepsilon}}
\]
where $\nabla u=du+\Gamma(dz)\cdot u$ is the Chern connection of
$(V\langle D\rangle,h_{V\langle D\rangle,\varepsilon})$, what we need to prove is that on the total space of $V$ over $X\setminus\abs{D}$, the $(1,1)$-form
\[
V\ni (z,u)\mapsto
\imath\ddbar\Abs{u}_{h_{V\langle D\rangle,\varepsilon}}^2+\gamma\,
\imath\ddbar\psi_\delta\,\Abs{u}_{h_{V\langle D\rangle,\varepsilon}},
\]
is non-negative. For this, we calculate the associated hermitian quadratic form on~$T_V$
\[
Q_{V\langle D\rangle,\gamma,\varepsilon}(z,u)(\xi,\eta),\quad (\xi,\eta)\in T_{V,(z,u)},\quad
\xi=\sum_{\ell=1}^n\xi_\ell\tfrac{\partial}{\partial z_\ell},\quad
\eta=\sum_{\lambda=1}^r\eta_\lambda\tfrac{\partial}{\partial u_\lambda},
\]
and observe that the curvature tensor is obtained by taking the restriction to the ``parallel'' directions $\nabla u=0$, that is, by substituting $du=-\Gamma(dz)\cdot u$, i.e.\
$\eta=-\Gamma(\xi)\cdot u$. Let us fix an arbitrary point
$z_0\in X\setminus\abs{D}$. We take local holomorphic coordinates $(z_1,\,\dotsc,\,z_n)$ centered at $z_0$, and let $(e_1,\,\dotsc,\,e_r)$ be a local holomorphic frame of $V$ such that
\[
\langle e_\lambda,e_\mu\rangle_{h_V}=\delta_{\lambda\mu}+
\sum_{\ell,m,\lambda,\mu}c_{\ell m\lambda\mu}\,z_\ell\bar z_m+O(\abs{z}^3),
\]
where the $\tfrac{\imath}{ 2\pi}c_{\ell m\lambda\mu}$ are the coefficients of
$-\Theta_{V,h_V}$. Let us write $u=\sum_{\lambda=1}^ru_\lambda e_\lambda$ and denote by $\langle u,v\rangle=\sum_{1\,\leq\,\lambda\,\leq\,r}u_\lambda\bar v_\lambda$ the standard hermitian form, $\abs{u}$ the associated norm. We find
\begin{multline*}
\tag{$\ref{eq:orbi_metric}_{0}$}\label{6.4_0}
\Abs{u}_{h_{V\langle D\rangle,\varepsilon}}^2 =
\abs{u}^2+\sum_{\ell,m,\lambda,\mu}c_{\ell m\lambda\mu}\,z_\ell\bar z_m u_\lambda\bar u_\mu+O(\abs{z}^3)
\\
+\sum_j\varepsilon_j\,\big(\abs{\sigma_j}^2e^{-\varphi_j}\big)^{-1+1/\rho_j}\,
\left|\partial\sigma_j(u)-\sigma_j\partial\varphi_j(u)\right|^2 e^{-\varphi_j},
\end{multline*}
since $\dbar\sigma_j=0$. In order to simplify the calculation, we set formally
\[
\left\{
\begin{aligned}
\tilde\sigma_j&=\sigma_j^{1/\rho_j}, &
\tilde\varepsilon_j&=\rho_j^2\varepsilon_j, &
\tilde\varphi_j&=\rho_j^{-1}\varphi_j, &&
\text{if }\rho_j<\infty, \\
\tilde\sigma_j&=\log\sigma_j, &
\tilde\varepsilon_j&=\varepsilon_j, &
\tilde\varphi_j&=\varphi_j,\hfill &&
\text{if }\rho_j=\infty.
\end{aligned}
\right.
\]
Respectively to the non-logarithmic and logarithmic situations, we then get the more tractable expression
\begin{multline}
\Abs{u}_{h_{V\langle D\rangle,\varepsilon}}^2 \\
=\abs{u}^2+{}\kern-4pt
\sum_{\ell,m,\lambda,\mu}\kern-4pt c_{\ell m\lambda\mu}\,z_\ell\bar z_m u_\lambda\bar u_\mu\,{+}\,O(\abs{z}^3)
\,{+}\sum_j\tilde\varepsilon_j\,
\left|\partial\tilde\sigma_j(u)\,{-}\,\sigma_j\partial\tilde\varphi_j(u)\right|^2 e^{-\tilde\varphi_j},
\end{multline}
\begin{multline}
\tag{$\theequation_{\infty}$}
\Abs{u}_{h_{V\langle D\rangle,\varepsilon}}^2 \\
=
\abs{u}^2+{}\kern-4pt
\sum_{\ell,m,\lambda,\mu}\kern-4pt c_{\ell m\lambda\mu}\,z_\ell\bar z_m u_\lambda\bar u_\mu\,{+}\,O(\abs{z}^3)
\,{+}\sum_j\tilde\varepsilon_j\,
\left|\partial\tilde\sigma_j(u)-\partial\tilde\varphi_j(u)\right|^2.
\end{multline}
More importantly, the poles have disappeared -- a fact reflecting the orbifold smoothness of the metric. In what follows, for the sake of simplicity, we remove the tildes in the notation, and conduct the calculation only in the non-logarithmic situation $(\rho_j<\infty)$, since the logarithmic case can be recovered by taking $\rho_j$ very large; this actually amounts to using a ramified change of variable $\tilde z'_\ell=\smash{z_\ell^{1/\rho_\ell}}$ in suitable coordinates, allowing us in this way to take $\rho_j=1$ in~\eqref{6.4_0}. Also, our later calculations will be done by adding the orbifold divisor components one by one. This essentially reduces the situation to the case where $D=(1-\tfrac{1}{\rho})\Delta$ only has one component, and the notation becomes much lighter. Therefore, we drop the indices $j$ and the summations $\sum_j$, and consider the simple situation where the metric is given by
\begin{align}\label{6.8}
\Abs{u}_{h_{V\langle D\rangle,\varepsilon}}^2 &=
\abs{u}^2+\sum_{\ell,m,\lambda,\mu}c_{\ell m\lambda\mu}\,z_\ell\bar z_m u_\lambda\bar u_\mu+O(\abs{z}^3)+\varepsilon\,|\partial\sigma(u) -\sigma\,\partial\varphi(u)|^2
\\
\tag{\theequation\ensuremath{'}}\label{6.8'}
\langle\!\langle u,v\rangle\!\rangle_{h_{V\langle D\rangle,\varepsilon}}^2 &=
\langle u,v\rangle^2+
\sum_{\ell,m,\lambda,\mu}c_{\ell m\lambda\mu}\,z_\ell\bar z_m u_\lambda\bar v_\mu+O(\abs{z}^3)
\\
\nonumber&
\qquad+\varepsilon\,\big(\partial\sigma(u)-\sigma\,\partial\varphi(u)\big)
\big(\,\overline{\partial\sigma(v)-\sigma\,\partial\varphi(v)}\,\big)
\,e^{-\varphi}.
\end{align}
We also take a holomorphic trivialization of the line bundle $\cO_X(\Delta)$ so that the associated weight $\varphi$ satisfies
$\varphi(z)=\sum_{\ell,m}\alpha_{\ell m}\,z_\ell\bar z_m+O(\abs{z}^3)$ near
$z_0=0$. Then
\[
\partial\varphi=\sum_{\ell,m}\alpha_{\ell m}\,\bar z_m dz_\ell+O(\abs{z}^2),
\quad
\dbar\varphi=\sum_{\ell,m}\alpha_{\ell m}\,z_\ell\,d\bar z_m+O(\abs{z}^2).
\]
At the point $z=z_0$, we have $\partial\varphi(z_0)=\partial\varphi(z_0)=0$,
$\nabla\sigma=\partial\sigma$, and our metric admits the expression
\[
\Abs{u}_{h_{V\langle D\rangle,\varepsilon}}^2=
\abs{u}^2+\varepsilon\,\abs{\partial\sigma(u)}^2,\quad
\langle\!\langle u,v\rangle\!\rangle_{h_{V\langle D\rangle,\varepsilon}}=
\langle u,v\rangle+\varepsilon\,\partial\sigma(u)\,
\overline{\partial\sigma(v)}.
\]
Let $u,v$ be arbitrary local holomorphic sections of $V$, and denote by
$\nabla_\xi$ the Chern covariant differentiation of $(V\langle D\rangle,h_{V\langle D\rangle,\varepsilon})$ in the direction
$\xi\in T_X$. By polarizing the quadratic form
$\Abs{u}_{h_{V\langle D\rangle,\varepsilon}}^2$ into a hermitian inner product
$\partial_\xi\langle\!\langle u,v
\rangle\!\rangle_{h_{V\langle D\rangle,\varepsilon}}$ and setting $\nabla_\xi u=
\nabla^{1,0}_\xi u=\partial_\xi u+\Gamma(\xi)\cdot u$, a differentiation of~\eqref{6.8'} at $z=z_0$ yields
\begin{align*}
\partial_\xi\langle\!\langle u,v
\rangle\!\rangle_{h_{V\langle D\rangle,\varepsilon}} &=
\langle\nabla_\xi u,v\rangle+\varepsilon\,
\partial\sigma(\nabla_\xi u)\,\overline{\partial\sigma(v)}
\\
&=
\langle\partial_\xi u\,,\,v\rangle+\varepsilon\;
\partial\sigma(\partial_\xi u)\;\overline{\partial\sigma(v)}
\\
&\phantom{=}\hspace{3\parindent} +\varepsilon\,\partial^2\sigma(\xi,u)\,\overline{\partial\sigma(v)} -\varepsilon\,
\partial\sigma(u)\,\bar \sigma\,\ddbar\varphi(\xi,v),
\end{align*}
where $\partial^2\sigma(\xi,u)\bydef\sum_\lambda\partial_\xi(\partial\sigma (e_\lambda)\,u_\lambda$ is viewed as an element of
$(T_X^*\otimes V^*)_{z_0}$ and $\ddbar\varphi$ as a hermitian form on $T_X$, operating on $T_X\otimes\bar V\subset T_X\otimes \bar T_X$. In fact,
$u\mapsto\partial\sigma(u)$ and $(\xi,u)\mapsto\partial^2\sigma(\xi,u)$ can be intrinsically defined as $\nabla^{1,0}\sigma_{|V}$ and
$\smash{\nabla_{V^*\otimes \cO(\Delta)}^{1,0}}(\nabla^{1,0}\sigma_{|V})$ at $z_0$, and we will denote them by $\nabla\sigma$ and $\nabla^2\sigma$. In this setting, a subtraction of the last two lines in our equalities shows that the $(1,0)$-form $\Gamma$ of the connection of $(V\langle D\rangle,h_{V\langle D\rangle})$ is given at $z_0$ by the formula
\begin{equation}\label{6.10}
\langle \Gamma(\xi)\cdot u,v\rangle +\varepsilon\,
\nabla\sigma(\Gamma(\xi)\cdot u)\,\overline{\nabla\sigma(v)} =\varepsilon\,
\nabla^2\sigma(\xi,u)\,\overline{\nabla\sigma(v)} -\varepsilon\,\nabla\sigma(u)\;\bar\sigma\,\ddbar\varphi(\xi,v).
\end{equation}
This equality is valid pointwise for any $u,v\in V_{z_0}$. As a consequence
\begin{equation}\label{6.11}
\Gamma(\xi)\cdot u+\varepsilon\,
\nabla\sigma(\Gamma(\xi)\cdot u)\,(\nabla\sigma)^* =\varepsilon\,
\nabla^2\sigma(\xi,u)\,(\nabla\sigma)^* -\varepsilon\,\nabla\sigma(u)\;\bar\sigma\, (\ddbar\varphi(\bullet,\xi))^*
\end{equation}

\break
\noindent where $\alpha^*\in V$ is the dual vector to a $1$-form $\alpha\in V^*$, such that $\langle\alpha^*,\cdot\rangle_{h_V}=\bar\alpha$. The special choice $v=\Gamma(\xi)\cdot u$ yields a (non-negative) real value in the left hand side of~\eqref{6.10}, and by taking the real part of the right hand side, we obtain
\stepcounter{equation}
\begin{multline}
\tag{$\theequation_{0}$}\label{6.12_0}
\abs{\Gamma(\xi)\cdot u}^2 +\varepsilon\,
\abs*{\nabla\sigma(\Gamma(\xi)\cdot u)}^2
\\
=\varepsilon\,\Re\big(
\nabla^2\sigma(\xi,u)\,
\overline{\nabla\sigma(\Gamma(\xi)\cdot u)}\,\big)-\varepsilon\,
\Re\big(
\nabla\sigma(u)\,\bar\sigma\;\ddbar\varphi(\xi,\Gamma(\xi)
\cdot u)\big).
\end{multline}
Also, by applying $\nabla\sigma$ to~\eqref{6.11}, we obtain
\begin{multline*}
\nabla\sigma(\Gamma(\xi)\cdot u)+\varepsilon\,
\nabla\sigma(\Gamma(\xi)\cdot u)\,
\langle\nabla\sigma,\nabla\sigma\rangle\\
=\varepsilon\,\nabla^2\sigma(\xi,u)\,
\left\langle\nabla\sigma,\nabla\sigma\right\rangle -\varepsilon\,\nabla\sigma(u)\;\bar\sigma\,
\left\langle\nabla\sigma,\ddbar\varphi(\bullet,\xi)\right\rangle,
\end{multline*}
hence
\begin{equation_}{1}\label{6.12_1}
\nabla\sigma(\Gamma(\xi)\cdot u)=
\tfrac{\varepsilon}{ 1+\varepsilon\,\abs{\nabla\sigma}^2}\big(
\nabla^2\sigma(\xi,u)\,\abs{\nabla\sigma}^2-\nabla\sigma(u)\;\bar\sigma\,
\left\langle\nabla\sigma,\ddbar\varphi(\bullet,\xi)\right\rangle\big).
\end{equation_}
As $2\pi\,\Theta_{A,h_A}=\imath\ddbar\psi_\delta$, we infer by a brute force calculation from~\eqref{6.8} that
\begin{align}
\nonumber\stepcounter{equation} Q_{V\langle D\rangle,\gamma,\varepsilon}(z,u)(\xi,\eta) &=
\ddbar\Abs{u}_{h_{V\langle D\rangle,\varepsilon}}^2\cdot(\xi,\eta) +\gamma\,\ddbar\psi_\delta(\xi,\xi)\,
\Abs{u}_{h_{V\langle D\rangle,\varepsilon}}^2
\\
\tag{$\theequation_1$}\label{6.13_1}
&=
\gamma\,\ddbar\psi_\delta(\xi,\xi)\,\abs{u}^2 +\sum_{\ell,m,\lambda,\mu}c_{\ell m\lambda\mu}\,\xi_\ell\bar\xi_m\, u_\lambda\bar u_\mu
\\
\tag{$\theequation_2$}\label{6.13_2}
&\qquad +\varepsilon\,\big(\gamma\,\ddbar\psi_\delta(\xi,\xi)-
\ddbar\varphi(\xi,\xi)\big)\,\abs{\nabla\sigma(u)}^2\kern80pt
\\
\tag{$\theequation_3$}\label{6.13_3}
&\qquad +\abs{\eta}^2+\varepsilon\,
\abs*{\nabla\sigma(\eta)+\nabla^2\sigma(\xi,u)}^2
\\
\tag{$\theequation_4$}\label{6.13_4}
&\qquad -2\varepsilon\,\Re\big(
\nabla\sigma(u)\,\bar\sigma\,\ddbar\varphi(\xi,\eta)\big)
\\
\tag{$\theequation_5$}\label{6.13_5}
&\qquad -2\varepsilon\,\Re\big(
\nabla\sigma(u)\,\ddbar\varphi(\xi,u)\,
\overline{\nabla\sigma(\xi)}\,\big)
\\
\tag{$\theequation_6$}\label{6.13_6}
&\qquad -2\varepsilon\,\Re\big(
\nabla\sigma(u)\,\bar\sigma\,\ddbar{\kern1pt}^2
\varphi(\xi,\xi,u)\big)
\\
\tag{$\theequation_7$}\label{6.13_7}
&\qquad +\varepsilon\,\abs{\sigma}^2\,\left|\ddbar\varphi(u,\xi)\right|^2,
\end{align}
where we identify a $(1,1)$-form such as $\ddbar\varphi$ with a hermitian form, and take $\eta=-\Gamma(\xi)\cdot u$. The second term in~\eqref{6.13_2} is obtained by differentiating
$\varepsilon\,\abs{\nabla\sigma(u)}^2$, while~\eqref{6.13_3}, \eqref{6.13_4} and~\eqref{6.13_5} actually come from the differentiation of the term $\dots\Re(\dots)$ in~\eqref{6.8}. By our assumptions~\eqref{6.3_j}, the first two terms~\eqref{6.13_1}, \eqref{6.13_2} are positive in the sense of Griffiths, and such that
\begin{align*}
\eqref{6.13_1}&\geq c\,\left(\abs{\xi}^2+\abs{\nabla\sigma(\xi)}^2\right)\,\abs{u}^2, \eqref{6.13_2}&\geq c\,\varepsilon\,\left(\abs{\xi}^2+\abs{\nabla\sigma(\xi)}^2\right)\,
\abs{\nabla\sigma(u)}^2,\quad c>0.
\end{align*}
(Here the term $\abs{\nabla\sigma(\xi)}^2$ is significant, because we will later replace $\sigma$ by $\sigma^{1/\rho}$ in the orbifold case, and then $\nabla\sigma^{1/\rho}(\xi)$ is unbounded with respect to $\abs{\xi})$. The third term~\eqref{6.13_3} is semi-positive. We claim that the terms (\ref{6.13_4}, \ref{6.13_5}, \ref{6.13_6}, \ref{6.13_7} are negligible for
$\varepsilon\ll 1$, in the sense that
$Q_{V\langle D\rangle,\gamma,\varepsilon}(z,u)(\xi,\eta)$ is comprised between $(1\pm\delta)$(\eqref{6.13_1}+\eqref{6.13_2}+\eqref{6.13_3}), with $\delta>0$ as small as we want when $\varepsilon\leq\varepsilon_0(\delta)$. In fact, since $\ddbar\varphi$ is smooth, there exists $C>0$ such that
\begin{align*}
\abs*{\eqref{6.13_4}} &\leq C\,\varepsilon\abs{\sigma}\,\abs{\nabla\sigma(u)}\,\abs{\xi}\,\abs{\eta}
\\
&\leq
\varepsilon^{3/2}\abs{\xi}^2
\abs{\nabla\sigma(u)}^2 + C^2\,\varepsilon^{1/2}
\abs{\sigma}^2
\abs{\eta}^2
\ll~\eqref{6.13_2}+\eqref{6.13_3}.
\end{align*}
Similarly
\begin{align*}
\abs*{\eqref{6.13_5}} &\leq C\,\varepsilon\, \abs{\xi}\,\abs{u}\,\abs{\nabla\sigma(\xi)}\,\abs{\nabla\sigma(u)}
\\
&\leq C\,\varepsilon^{3/2}\,\abs{\xi}^2\,\abs{\nabla\sigma(u)}^2 +C\,\varepsilon^{1/2}\,\abs{\nabla\sigma(\xi)}^2\,\abs{u}^2
\ll~\eqref{6.13_1}+\eqref{6.13_2}.
\end{align*}
The last two terms (\eqref{6.13_6}, \eqref{6.13_7}) are even easier, since
\begin{align*}
\abs*{\eqref{6.13_6}} &\leq C\,\varepsilon\,\abs{\sigma}\,\abs{\xi}^2\,\abs{u}\,\abs{\nabla\sigma(u)}
\leq \varepsilon^{1/2}\abs{\xi}^2\,\abs{u}^2+ C^2\,\varepsilon^{3/2}\,\abs{\sigma}^2\,\abs{\xi}^2\,\abs{\nabla\sigma(u)}^2
\\
&\ll~\eqref{6.13_1} +~\eqref{6.13_2},
\\[4pt]
\abs*{\eqref{6.13_7}} &\leq C\,\varepsilon\,\abs{\xi}^2\abs{u}^2
\ll~\eqref{6.13_1}.
\end{align*}
Finally, by replacing $\eta$ with $-\Gamma(\xi)\cdot u$ and using (\eqref{6.12_0}, \eqref{6.12_1}), we find
\begin{align}
\nonumber\stepcounter{equation}~\eqref{6.13_3}+\eqref{6.13_4} &=
\abs*{\Gamma(\xi)\cdot u}^2\\ &\quad{}+\varepsilon\,
\big|\nabla\sigma(\Gamma(\xi)\,{\cdot}\,u) -\nabla^2\sigma(\xi,u)\big|^2 +2\varepsilon\,\Re\big(
\nabla\sigma(u)\,\bar\sigma\,
\ddbar\varphi(\xi,\Gamma(\xi)\cdot u)\big)\\ &=\eqref{6.12_0}+\varepsilon\,\big|
\nabla^2\sigma(\xi,u)\big|^2-2\varepsilon\,\Re\big(
\,\overline{\nabla^2\sigma(\xi,u)}\,
\nabla\sigma(\Gamma(\xi)\cdot u)\big)\\ &\quad{}+2\varepsilon\,\Re\big(
\nabla\sigma(u)\,\bar\sigma\,
\ddbar\varphi(\xi,\Gamma(\xi)\cdot u)\big)\\ &=
\varepsilon
\abs*{\nabla^2\sigma(\xi,u)}^2 -\varepsilon\,\Re\big(\,
\overline{\nabla^2\sigma(\xi,u)}\,
\nabla\sigma(\Gamma(\xi)\cdot u)\big)\\ &\quad{}+\varepsilon\,\Re\big(
\nabla\sigma(u)\,\bar\sigma\,
\ddbar\varphi(\xi,\Gamma(\xi)\cdot u)\big).\\
\tag{$\theequation_1$}\label{6.14_1}
&=\tfrac{\varepsilon}{ 1+\varepsilon\abs{\nabla\sigma}^2}\,
\abs*{\nabla^2\sigma(\xi,u)}^2
\\
\tag{$\theequation_2$}\label{6.14_2}
&\quad{}+\tfrac{\varepsilon}{ 1+\varepsilon\,\abs{\nabla\sigma}^2}\,\Re\Big(
\overline{\nabla^2\sigma(\xi,u)}\,
\varepsilon\,\nabla\sigma(u)\,\bar\sigma\,
\left\langle\nabla\sigma,\ddbar\varphi(\bullet,\xi)\right\rangle\Big)
\\
\tag{$\theequation_3$}\label{6.14_3}
&\quad{}+\varepsilon\,\Re\big(
\nabla\sigma(u)\,\bar\sigma\,
\ddbar\varphi(\xi,\Gamma(\xi)\cdot u)\big).
\end{align}
The term~\eqref{6.14_3} equals $\tfrac{1}{2}\eqref{6.13_4}$, thus it is negligible, and the term~\eqref{6.14_2} admits an obvious bound
\begin{align*}
~\eqref{6.14_2} &\leq
\tfrac{\varepsilon}{ 1+\varepsilon\abs{\nabla\sigma}^2}
\Big(\varepsilon^{1/2}\,\left|\nabla^2\sigma(\xi,u)\right|^2+\varepsilon^{3/2}\,
\abs{\sigma}^2\,\abs{\nabla\sigma}^2\,\abs{\nabla\sigma(u)}^2\,\abs{\xi}^2\Big)\\
&\leq
\varepsilon^{1/2}\eqref{6.14_1}+
\varepsilon^{3/2}\abs{\sigma}^2\,\abs{\nabla\sigma(u)}^2\,\abs{\xi}^2\ll~\eqref{6.14_1}+\eqref{6.13_2}.
\end{align*}
By collecting all non-negligible terms~\eqref{6.13_1}, \eqref{6.13_2} and~\eqref{6.14_1}, we obtain a curvature form
\begin{multline*}
Q_{V\langle D\rangle,\gamma,\varepsilon}(z)(\xi\otimes u)
\simeq{}\gamma\,\ddbar\psi_\delta(\xi,\xi)\,\abs{u}^2 +\sum_{\ell,m,\lambda,\mu}c_{\ell m\lambda\mu}\,\xi_\ell\bar\xi_m\, u_\lambda\bar u_\mu\\
+\varepsilon\,
\big(\gamma\,\ddbar\psi_\delta(\xi,\xi)-\ddbar\varphi(\xi,\xi)\big)
\,\abs{\nabla\sigma(u)}^2 +\tfrac{\varepsilon}{ 1+\varepsilon\abs{\nabla\sigma}^2}\,
\abs*{\nabla^2\sigma(\xi,u)}^2.
\end{multline*}
At this point, we come back to the orbifold situation, and thus replace $\sigma$ by $\sigma^{1/\rho}$,
$\varphi$ by $\rho^{-1}\varphi$ and $\varepsilon$ by
$\rho^2\,\varepsilon$. This gives the curvature estimate
\begin{multline}
Q_{V\langle D\rangle,\gamma,\varepsilon}(z)(\xi\otimes u)\simeq{}
\gamma\,\ddbar\psi_\delta(\xi,\xi)\,\abs{u}^2 +\sum_{\ell,m,\lambda,\mu}c_{\ell m\lambda\mu}\,\xi_\ell\bar\xi_m\, u_\lambda\bar u_\mu\\
\begin{aligned}
&+\varepsilon\,\abs{\sigma}^{-2+2/\rho}\,
\big(\gamma\,\ddbar\psi_\delta(\xi,\xi)-\rho^{-1}\,\ddbar\varphi(\xi,\xi)\big)
\,\abs{\nabla\sigma(u)}^2
\\
&+\tfrac{\varepsilon\,\abs{\sigma}^{-2+2/\rho}}{ 1+\varepsilon\,\abs{\sigma}^{-2+2/\rho}\,
\abs{\nabla\sigma}^2}\,\big|\,\nabla^2\sigma(\xi,u) -(1-1/\rho)\,\sigma^{-1}\,\nabla\sigma(\xi)\nabla\sigma(u)\big|^2,
\end{aligned}
\end{multline}
In the general situation $D=\sum_{1\,\leq\,j\,\leq\,N}(1-1/\rho_j)\Delta_j$ of a multi-component orbifold divisor, we add the components $\Delta_j$ one by one, and obtain inductively the following quantitative estimate, which is a rephrasing of Theorem~\ref{theo:orbi_curv_estimate}.
\end{proof}

\end{document}
